\section{Methods}
\label{sec:methods}

\subsection{Compiled representation and composition}
\label{sec:d2l}

We use the released D2L compilers~\citep{charakorn2026doctolora} for
Gemma-2-2B-it~\citep{gemmateam2024gemma2} and Qwen3-4B-Instruct-2507~\citep{yang2025qwen3}.
The architecture, composition rule and bias adapter described here are from D2L.
We introduce the per-rank notation and gating interventions below.
Figure~\ref{fig:architecture} shows the native D2L pipeline. At write time, a
frozen copy of the target model reads a document $d$. A Perceiver-style
hypernetwork then converts the resulting layer activations into a LoRA
adapter: for each MLP layer $l$, it produces two rank-8 matrices,
$A_{l,d}$ and $B_{l,d}$, for the layer's down-projection. The hypernetwork was
trained by context distillation. The target model is frozen; compiler and
adapter updates are restricted to the fine-tuning and passage-fitting controls.
Thus, a document is compiled once
into model state and does not need to remain in the serving prompt.

We describe one layer at one token position, omitting layer and token indices
throughout the equations below. Let $x$ be the MLP down-projection input,
$z$ its base output, and $z'$ its modified output. The residual $h_{\mathrm{res}}$
enters the MLP branch after attention, and $h_{\mathrm{out}}$ is the block output
passed to the next layer. $\mathcal{N}$ is Gemma's post-MLP normalization,
or the identity for Qwen.

For document $d$, $A_d$ maps $x$ into rank space and $B_d$ maps back to the
down-projection output space. We include all adapter scaling in these factors,
absorbing the fixed implementation scale into $B_d$. The same convention
applies to the shared bias factors $A_0,B_0$. Native D2L sums over the supplied documents:
\begin{equation}
\begin{aligned}
z&=W_{\mathrm{down}}x\\
z'&=z+\sum_d B_dA_dx+B_0A_0x\\
h_{\mathrm{out}}&=h_{\mathrm{res}}+\mathcal{N}(z')
\end{aligned}
\label{eq:sum}
\end{equation}
Both document and bias adapters have rank $R=8$. The shared bias is added once;
the subscript 0 does not denote the first document or rank.
Appendix~\ref{sec:extended-methods} relates this matrix orientation to the
row-stored factors in the implementation.

\begin{figure*}[t]
\centering
\resizebox{0.92\textwidth}{!}{%
\begin{tikzpicture}[font=\footnotesize,>={Stealth[length=4pt]},
 box/.style={draw,rounded corners=2pt,align=center,inner sep=3pt,line width=0.5pt},
 input/.style={box,fill=orange!12,draw=orange!60!black},
 frozen/.style={box,fill=gray!10,draw=gray!70},
 trained/.style={box,fill=violet!10,draw=violet!60},
 gen/.style={box,fill=teal!12,draw=teal!70!black},
 flow/.style={->,draw=gray!70!black,line width=0.6pt},
 tap/.style={->,draw=teal!60!black,line width=0.6pt},
 rowlabel/.style={anchor=west,font=\footnotesize\bfseries,text=gray!70!black}]
\def\xa{3.6}\def\xb{7.3}\def\xc{11.9}
% ---------------- write time ----------------
\node[rowlabel] at (-1.2,1.05) {Write time (once per document)};
\node[input,minimum width=2.2cm,minimum height=1.0cm] (d3) at (0.24,-0.16) {};
\node[input,minimum width=2.2cm,minimum height=1.0cm] (d2) at (0.12,-0.08) {};
\node[input,minimum width=2.2cm,minimum height=1.0cm] (doc) at (0,0) {Document $d$};
\node[anchor=north,align=center,font=\scriptsize] at (0.1,-0.75)
  {bank: source $d_1$,\\distractors $d_2,\dots,d_K$};
\foreach \x/\n/\t in {\xa/w0/Layer $0$,\xb/w1/Layer $1$,\xc/wl/Layer $l$}
  \node[frozen,minimum width=2.3cm,minimum height=0.8cm] (\n) at (\x,0) {\t};
\node at ($(w1)!0.5!(wl)$) {$\cdots$};
\node (wend) at (\xc+2.0,0) {$\cdots$};
\draw[flow] (d3.east |- doc) -- (w0);
\draw[flow] (w0) -- (w1);
\draw[flow] (w1.east) -- ++(0.6,0);
\draw[flow] ($(w1)!0.5!(wl)+(0.4,0)$) -- (wl);
\draw[flow] (wl) -- (wend);
% ---------------- compiler ----------------
\foreach \x/\n/\i in {\xa/h0/0,\xb/h1/1,\xc/hl/l}
  \node[trained,minimum width=2.3cm,minimum height=0.85cm] (\n) at (\x,-2.0)
    {Perceiver $\psi$};
\foreach \x/\w/\h/\i in {\xa/w0/h0/0,\xb/w1/h1/1,\xc/wl/hl/l}{
  \draw[tap] (\x-1.55,0) -- (\x-1.55,-1.0) -- (\x,-1.0) -- (\h.north);
  \node[anchor=west,font=\scriptsize] at (\x+0.05,-0.95) {layer-$\i$ input};
}
% ---------------- generated factors ----------------
\foreach \x/\n/\i in {\xa/a0/0,\xb/a1/1,\xc/al/l}
  \node[gen,minimum width=2.3cm,minimum height=0.75cm] (\n) at (\x,-3.45)
    {LoRA};
\foreach \h/\a in {h0/a0,h1/a1,hl/al} \draw[tap] (\h) -- (\a);
\node[anchor=west,align=left,font=\scriptsize] at (\xc+1.35,-3.45)
  {For a bank, the factors of\\$d_1,\dots,d_K$ are concatenated\\along the rank dimension:\\$8K$ components, unit weight};
% ---------------- serving time ----------------
\node[rowlabel] at (-1.2,-4.55) {Serving time (every request)};
\node[input,minimum width=2.2cm,minimum height=0.9cm] (q) at (0,-5.7) {Question\\(no document text)};
\foreach \x/\n/\t in {\xa/s0/Layer $0$,\xb/s1/Layer $1$,\xc/sl/Layer $l$}{
  \node[frozen,minimum width=3.15cm,minimum height=2.0cm] (\n) at (\x,-6.0) {};
  \node[font=\footnotesize] at (\x,-5.2) {\t};
  \node[box,fill=white,font=\scriptsize,inner sep=2pt] (\n-attn) at (\x-.85,-5.75) {Attention};
  \node[box,fill=white,font=\scriptsize,inner sep=2pt] (\n-mlp) at (\x+.35,-5.75) {MLP\\up/gate};
  \node[box,fill=teal!8,draw=teal!70!black,font=\scriptsize,inner sep=2pt]
    (\n-down) at (\x+.35,-6.45) {Down-proj.};
  \node[draw=teal!70!black,circle,fill=white,inner sep=1pt,font=\scriptsize]
    (\n-add) at (\x+1.25,-6.45) {$+$};
  \draw[flow] (\n-attn) -- (\n-mlp);
  \draw[flow] (\n-mlp) -- (\n-down);
  \draw[flow] (\n-down) -- (\n-add);
}
\node (gap) at (9.55,-5.75) {$\cdots$};
\node (send) at (\xc+2.1,-6.45) {$\cdots$};
\node[box,minimum width=1.5cm,minimum height=0.9cm] (ans) at (\xc+3.6,-6.45) {Answer};
\draw[flow] (q) -- (s0-attn);
\draw[flow] (s0-add) -- ++(.5,0) |- (s1-attn);
\draw[flow] (s1-add) -- ++(.5,0) |- (gap);
\draw[flow] (gap) -- (sl-attn);
\draw[flow] (sl-add) -- (send);
\draw[flow] (send) -- (ans);
\foreach \x/\a/\s in {\xa/a0/s0,\xb/a1/s1,\xc/al/sl}
  \draw[tap] (\a.south) -- (\x,-4.75) -- (\x+1.25,-4.75) -- (\s-add.north);
% ---------------- legend ----------------
\begin{scope}[shift={(-1.15,-7.65)},font=\scriptsize]
\node[input,minimum width=0.35cm,minimum height=0.25cm,inner sep=0pt] (k1) at (0,0) {};
\node[anchor=west] at (0.25,0) {input text};
\node[frozen,minimum width=0.35cm,minimum height=0.25cm,inner sep=0pt] at (2.0,0) {};
\node[anchor=west] at (2.25,0) {frozen model};
\node[trained,minimum width=0.35cm,minimum height=0.25cm,inner sep=0pt] at (4.3,0) {};
\node[anchor=west] at (4.55,0) {trained compiler};
\node[gen,minimum width=0.35cm,minimum height=0.25cm,inner sep=0pt] at (6.9,0) {};
\node[anchor=west] at (7.15,0) {generated adapter};
\end{scope}
\end{tikzpicture}}
\caption{Doc-to-LoRA. At write time the frozen model reads a document, and the
Perceiver compiler $\psi$ turns the input to each layer $l$ into rank-8 factors $A_{l,d}$ and
$B_{l,d}$. The retention probe reads the latent between the Perceiver and the output
head. One compiler serves all layers, with a separate output head for each layer. At serving time the prompt
holds only the question. The teal arrows show where LoRA adds
the document and shared bias contributions to the base down-projection output
$z$, producing the modified output $z'$ (Equation~\ref{eq:sum}).
The student-layer detail omits normalization and residual connections.
Each document in a bank is compiled
separately, and the factors of the source and the distractors are concatenated along
the rank dimension.}
\label{fig:architecture}
\end{figure*}


The ICL reference gives the same model the same documents as text in its
context. We call this text-conditioned model the \emph{teacher}, because its
behavior is the reference that the compiled state is intended to reproduce. A
single-document comparison measures how closely the adapter reproduces the
model's behavior when it reads the document directly. For composition, a \emph{bank} is
the set of $K$ documents whose adapter contributions are summed for one request
(Figure~\ref{fig:architecture}). Each bank has one answer-bearing document,
called the \emph{source}, and $K-1$ documents that do not answer the question,
called \emph{distractors}. This arrangement tests whether the source's useful
information survives the addition of unrelated compiled states.
This is a distractor-robustness test, not a test of combining jointly necessary
sources. Source selection is an oracle control here; a general user-memory
system must also allow several memories to contribute to one request.

\subsection{Evaluation sets}
\label{sec:suites}

We use SQuAD, DROP, and ROPES, the factual-QA benchmarks used by the original
D2L study~\citep{rajpurkar2016squad,dua2019drop,lin2019ropes,charakorn2026doctolora}.
Unless a paragraph states otherwise, the metric is token F1. Qwen answers in full sentences, which
token F1 scores poorly against short reference spans, so for Qwen we also report
answer containment and compare each condition with Qwen's own single-document
score. Appendix~\ref{sec:set-details} gives the
construction and size of every set.

\begin{table}[t]
\centering
\caption{Factual-QA evaluation sets. Diagnostic baselines are specific to
the set and intervention runtime; their scores are not pooled.}
\label{tab:evaluation-sets}
\scriptsize\setlength{\tabcolsep}{3pt}
\begin{tabular}{@{}lrl@{}}
\toprule
Set & Questions & Purpose / results \\
\midrule
SQuAD test (200 sources) & 1,127 & Reader comparison, Fig.~\ref{fig:composition} \\
SQuAD development (180) & 1,260 & Localization, Fig.~\ref{fig:layers} \\
SQuAD diagnostic (57) & 400 & Readout, Table~\ref{tab:readout-summary} \\
Mixed-task diagnostic & 402 & Transplants/masks, Table~\ref{tab:causal} \\
SQuAD passage fits (160) & 864 & Rank/coverage, Appendix~\ref{sec:single-tables} \\
DROP/ROPES (612) & 1,934 & Localization, Fig.~\ref{fig:layers}(b) \\
\bottomrule
\end{tabular}
\end{table}

\subsection{Single-document diagnostics}
\label{sec:diagnostics}

\textbf{Retention (RQ1.1).} A trained probe predicts passage-token presence
from the compiler latent before conversion into LoRA factors. This measures
the latent, not the final adapter. Separately, we check whether incorrect D2L
answers match passage spans or another question's reference answer.

\textbf{Readout, compiler, and rank (RQ1.2).} We compare the post-block,
write-site, and gated readers defined in Section~\ref{sec:readers} with native D2L, fine-tune the compiler alone or
jointly with a reader, and fit rank-8--64 adapters to individual passages using
generated questions. We report passage-fit accuracy separately for answers
covered and not covered by those questions. We also test a question-conditioned compiler.

\textbf{Missing computation (RQ1.3).} Hidden-state transplants replace native
states with teacher states at question positions, question plus answer-history
positions, or all shared positions. Document-attention masks restrict which
positions can read the document.

\subsection{Multi-document diagnostics}
\label{sec:composition-diagnostics}

\textbf{Update size (RQ2.1).} Uniform normalization scales each document's
$A_d$ by $K^{-\beta}$. Development data select $\beta=0.5$ on both models.
This reduces update magnitude without selecting documents.

A \emph{source-only mask} retains the known source's adapter contribution and
removes all distractor contributions. A \emph{silent bank} removes all document
contributions. Both retain the shared bias. These diagnostic masks are prescribed;
Section~\ref{sec:readers} describes learned readers, including predicted rank gates.

\textbf{Damage and useful work by layer (RQ2.2).} A \emph{layer-prefix oracle}
retains native composition below layer $L$ and applies the source-only mask
from $L$ onward, measuring interference accumulated in the earlier layers.
A \emph{bank-silent schedule} instead silences the bank below $L$, measuring
the cost of withholding early source contributions. Additional interventions
delay masking or source application until the end of prompt processing,
testing whether later intervention can recover accuracy.
Appendix~\ref{sec:composition-details} specifies the schedules and cache handling.

\textbf{Separability (RQ2.3).} Membership probes predict source membership
from a layer's hidden state and a key for each rank. We measure held-out AUC
and test whether gates derived from the probes improve generated answers.

\subsection{Trained readers}
\label{sec:readers}

We test three readers that modify how frozen document adapters affect the model,
targeting single-document accuracy (RQ1.2) and distractor interference (RQ2.3).

In reader inputs, $A_d,B_d$ denote the factors for all documents in the bank.
A \emph{gated reader} predicts a vector $g_d$ of eight rank multipliers
for each document:
\begin{equation}
\begin{aligned}
g_d&=\boldsymbol{G_d}(h_{\mathrm{res}},x;A_d,B_d)\\
z'&=z+\sum_d B_d(\boldsymbol{g_d\odot} A_dx)+B_0A_0x
\end{aligned}
\label{eq:gated-reader}
\end{equation}
A \emph{write-site reader} uses the native down-projection output to predict
a correction, then adds it to that output:
\begin{equation}
z'\leftarrow z'+\boldsymbol{R_{\mathrm{write}}}(h_{\mathrm{res}},x,z';A_d,B_d)
\label{eq:write-site-reader}
\end{equation}
A \emph{post-block reader} instead adds its correction to the native block
output. It receives the individual rank effects and the shared bias effect:
\begin{equation}
\begin{aligned}
E_d&=B_d\operatorname{diag}(A_dx)\\
h'_{\mathrm{out}}&=h_{\mathrm{out}}+\boldsymbol{R_{\mathrm{post}}}(h_{\mathrm{out}};E_d,B_0A_0x)
\end{aligned}
\label{eq:post-block-reader}
\end{equation}
Here $\operatorname{diag}$ places a vector's entries on a diagonal matrix.
Each column of $E_d$ is one rank's contribution; the reader receives $E_d$
for every document.
This correction is added after normalization and residual addition.

The compiled factors and shared bias remain frozen. All readers are trained on
SQuAD; the single-document reference uses $K=1$ only, while the composition
readers use mixed-$K$ banks. Gates start at one and additive corrections at
zero, so every reader initially reproduces native D2L. Gated readers also
receive source-membership supervision. Appendix~\ref{sec:rank-gates} gives the
architectures and training protocols.

\subsection{D2L on MidBrain}
\label{sec:deployment}

\textbf{MidBrain on PersonaMem-v2 (RQ3).} The memories are produced by the
unchanged production MidBrain pipeline. It combines episodic and consolidated
retrieval, retaining episodic memories as unmodified source material.
Individual episodes are expected to be accessed infrequently, whereas the
smaller set of consolidated memories is reused more often. We therefore
compile consolidated memory, where repeated use can amortize compilation,
and retain retrieved episodes as text.

We use full history as text as the
baseline and compare MidBrain's
consolidated-memory prompt supplied as text with the same prompt compiled
jointly as one document per user. This includes persona/profile context and
consolidated entries. Retrieved episodes remain text, with the same episodes
supplied in both production conditions. We also compare these conditions
without episodic retrieval and test compiling consolidated entries separately.
For the matched multi-adapter comparison, we compile persona/profile context
as one document and each consolidated entry separately, retaining the same
retrieved episodes as text. We compare native summation with $K^{-0.5}$
normalization, where $K$ includes the persona/profile adapter.
We score the four option labels directly at the first answer position,
without generating chain-of-thought reasoning. This differs from the original
PersonaMem-v2 evaluation, which permits reasoning before the final
answer~\citep{jiang2025personamemv2}, and may contribute to differences in
absolute accuracy.

\textbf{Cost (RQ3).} We compare request-time prefill and decoding with
document text in context against prefill and decoding with its compiled
LoRA adapter. We calculate per-request MACs, attended positions, and
resident bytes for both conditions.
Separately, we measure latency, throughput, and GPU memory for ICL with full prefill and D2L on one
A100 using the released implementation. Throughput is measured with fixed
batch sizes of 1, 8, and 16, with one request and distinct memory per user;
this measures batched execution rather than concurrent-request scheduling.
Compilation is a write-time cost excluded from the request-time comparison.
